% =====================================================================
% La otra mitad del dato: los precios
% Parte I — El dato
%
% ENCARGO: bolsa XM, techo de escasez, tarifa CEDENAR y su desglose,
% LCOE solar y la banda de precios que resulta.
%
% Reglas del documento:
% - Una idea = una subseccion = una figura.
% - Toda figura declara su procedencia con \fuente{...}.
% - Toda cifra sale del canon de agosto; ver GUIA.md.
% - M1 y M3 se reportan siempre en paralelo.
% - Ningun superlativo sin releer el CSV de la frontera correcta.
% =====================================================================
\section{La otra mitad del dato: los precios}
\label{sec:precios}
\justifying

Los tres capítulos anteriores construyeron una pareja de matrices de
energía. Con ellas todavía no se puede decir nada económico, porque un
kilovatio-hora no tiene un valor único: vale una cosa si la institución lo
compra a su comercializador, otra si lo inyecta a la red, y otra distinta
si se lo compra a la institución vecina. Este capítulo reconstruye esa
segunda mitad del dato desde las fuentes primarias, y termina delimitando
la banda dentro de la cual puede existir un intercambio que convenga a las
dos partes.

Conviene anticipar por qué el capítulo es más largo de lo que parecería
necesario. Todas las cifras del Capítulo~\ref{sec:res} son productos de
una energía por un precio. La energía se midió; el precio se construyó. Un
error de un diez por ciento en la construcción del precio es un error de
un diez por ciento en el resultado, y por eso cada componente del precio
se documenta con su fuente.

\subsection{El precio al que la red compra}
\label{sub:precios-bolsa}

Cuando una institución produce más de lo que consume y no encuentra
comprador dentro de la comunidad, el excedente va a la red y se liquida al
precio del mercado mayorista. Esa serie es, por tanto, el suelo económico
de todo el ejercicio: ningún vendedor aceptará dentro de la comunidad un
precio inferior al que obtendría fuera. La
Figura~\ref{fig:precios-bolsa} presenta la serie completa junto a su
patrón intradiario.

\begin{figure}[H]
 \centering
 \includegraphics[width=0.95\textwidth]{figuras/f5_01_bolsa.png}
 \caption[Serie horaria del precio de bolsa]{Precio del mercado
 mayorista sobre el horizonte del estudio. A la izquierda, la serie
 horaria en gris con la media diaria superpuesta; a la derecha, el
 patrón intradiario con su recorrido intercuartílico.
 \fuente{elaboración propia desde
 \file{data/precios\_bolsa\_xm\_api.csv}, recortado a las \num{6144}
 horas del horizonte.}}
 \label{fig:precios-bolsa}
\end{figure}

\notafig{La media del horizonte es de \uni{182.5}{COP/kWh} y la mediana
de \uni{114.3}{COP/kWh}. La distancia entre ambas no es un detalle
estadístico: indica una serie muy asimétrica, con la masa concentrada en
valores bajos y una cola larga de horas caras que llega a
\uni{2224.0}{COP/kWh}. Elaboración propia desde
\file{data/precios\_bolsa\_xm\_api.csv}.}

\begin{trampabox}
El archivo de caché de precios cubre hasta enero de 2026, mientras que el
horizonte del estudio termina el 16 de diciembre de 2025. Promediar el
archivo completo da \uni{193.6}{COP/kWh} en lugar de \uni{182.5}{}, una
diferencia del \pct{6.1} que se propagaría a toda comparación con la
tarifa. Todas las figuras y cifras de este capítulo usan la serie
recortada al horizonte; el recorte lo aplica el propio lector de datos,
para que no dependa de que cada script lo recuerde.
\end{trampabox}

\subsection{El precio al que la red vende}
\label{sub:precios-cu}

En el otro extremo está lo que la institución paga por la energía que
compra. En Colombia ese precio es el costo unitario, una suma de siete
componentes fijadas por resolución, y su composición importa aquí más de
lo que suele importar en un estudio energético: cada uno de los mecanismos
regulatorios del Capítulo~\ref{sec:esc} exime al usuario de componentes
distintas. Sin el desglose, esas diferencias no se pueden ni enunciar. La
Figura~\ref{fig:precios-cu} descompone el costo unitario mes a mes y
ordena sus componentes por peso.

\begin{figure}[H]
 \centering
 \includegraphics[width=0.95\textwidth]{figuras/f5_03_cu_desglose.png}
 \caption[Composición del costo unitario]{A la izquierda, el costo
 unitario mes a mes desagregado en sus siete componentes, para el nivel
 de tensión 2 en categoría oficial. A la derecha, el peso medio de cada
 componente sobre el horizonte.
 \fuente{elaboración propia desde
 \file{data/tarifas\_cedenar\_mensual.csv}, transcrito de los boletines
 mensuales del comercializador archivados en
 \file{data/cedenar\_pdfs/}.}}
 \label{fig:precios-cu}
\end{figure}

\notafig{El costo unitario medio del horizonte es de
\uni{792.06}{COP/kWh} en categoría oficial, con un recorrido anual de
\uni{50.18}{COP/kWh}, mínimo en enero de 2026 y máximo en agosto de 2025.
La estabilidad de la barra apilada a lo largo de los trece meses indica
que ninguna componente sufre saltos regulatorios dentro del horizonte.
Elaboración propia desde \file{data/tarifas\_cedenar\_mensual.csv}.}

Su reparto medio es:
generación, \uni{305.6}{COP/kWh} (\pct{38.6}); comercialización,
\uni{174.6}{} (\pct{22.0}); distribución, \uni{157.6}{} (\pct{19.9});
transmisión, \uni{54.9}{} (\pct{6.9}); otros cargos, \uni{41.2}{}
(\pct{5.2}); restricciones, \uni{36.8}{} (\pct{4.6}); y pérdidas,
\uni{21.3}{} (\pct{2.7}).

\begin{lecturabox}
La cifra que conviene retener de esta figura es que \textbf{la generación
apenas explica el \pct{38.6} de lo que se paga}. El resto, casi dos
tercios, son cargos de red, comercialización y conceptos regulatorios. Ahí
está la razón económica de fondo de todo el trabajo: un intercambio que
ocurre dentro de la comunidad, sin usar la red de distribución, puede
evitar una fracción de esos dos tercios, y esa fracción evitada es mayor
que cualquier ganancia de eficiencia en la generación misma. El
Capítulo~\ref{sec:disc} vuelve sobre este punto porque es también la
objeción más seria que se le puede hacer al mecanismo.
\end{lecturabox}

\subsection{Dos tarifas para la misma comunidad}
\label{sub:precios-categorias}

Las cinco instituciones no pagan lo mismo. Dos de ellas, Udenar y el
HUDN, son entidades públicas y pagan el costo unitario exacto; las otras
tres (Mariana, la UCC y Cesmag) son privadas y pagan además la
contribución de solidaridad del veinte por ciento. La
Figura~\ref{fig:precios-categorias} pone las dos series sobre un mismo eje
y sombrea la distancia entre ellas.

\begin{figure}[H]
 \centering
 \includegraphics[width=0.95\textwidth]{figuras/f5_04_oficial_comercial.png}
 \caption[Costo unitario por categoría tarifaria]{Costo unitario
 mensual de las dos categorías presentes en la comunidad. La franja
 sombreada entre ambas es la contribución de solidaridad.
 \fuente{elaboración propia desde
 \file{data/tarifas\_cedenar\_mensual.csv}.}}
 \label{fig:precios-categorias}
\end{figure}

\notafig{La diferencia media es de \uni{158.4}{COP/kWh}, exactamente el
\pct{20.0} del costo oficial. La razón entre las dos series vale
$1{,}2000$ en los trece meses, sin una sola excepción, lo que confirma que
se trata de un factor regulatorio y no de una diferencia de condiciones de
suministro. Elaboración propia desde
\file{data/tarifas\_cedenar\_mensual.csv}.}

Esta heterogeneidad tiene una consecuencia de diseño que conviene dejar
explícita, porque es una de las decisiones metodológicas menos evidentes
del trabajo: en la liquidación de los escenarios regulatorios, el precio
al que cada institución compra no es un número sino una matriz, con una
fila por institución y una columna por hora, que asigna a cada hora el
costo unitario vigente para esa institución en el mes que la contiene. Una
comunidad de cinco miembros con dos categorías tarifarias y trece meses de
tarifas distintas no puede liquidarse con un escalar sin introducir un
sesgo sistemático.

\begin{detallebox}
El juego del Capítulo~\ref{sec:modelo}, en cambio, sí usa un escalar
comunitario. La asimetría es deliberada y conviene entenderla: la
dinámica que busca el precio de equilibrio necesita un único techo común
(si cada agente tuviera el suyo, no habría un precio de mercado sino
cinco), mientras que la liquidación posterior sí puede y debe respetar la
tarifa de cada uno. El mecanismo que reconcilia ambas cosas es un techo
por agente aplicado en el momento de liquidar, que el
Capítulo~\ref{sec:modelo} desarrolla.
\end{detallebox}

\subsection{La banda de ganancia mutua}
\label{sub:precios-banda}

Con las dos series construidas queda delimitado el espacio dentro del cual
un intercambio entre pares tiene sentido económico. La condición es doble
y elemental: el precio debe superar lo que el vendedor obtendría
inyectando a la red, y quedar por debajo de lo que el comprador pagaría
comprándola a su comercializador. Si se cumple, las dos partes ganan; si
no se cumple, al menos una prefiere la red. La
Figura~\ref{fig:precios-banda} reúne los cuatro niveles construidos en
este capítulo y delimita con ellos la banda resultante.

\begin{figure}[H]
 \centering
 \includegraphics[width=0.95\textwidth]{figuras/f5_07_banda.png}
 \caption[Piso y techo del precio entre pares]{Los cuatro niveles de
 referencia construidos en este capítulo y la banda que delimitan. Todo
 precio de equilibrio del mercado entre pares vive dentro de la franja
 sombreada.
 \fuente{elaboración propia desde
 \file{data/precios\_bolsa\_xm\_api.csv} y
 \file{data/tarifas\_cedenar\_mensual.csv}.}}
 \label{fig:precios-banda}
\end{figure}

\notafig{La banda no procede de una diferencia de costos de producción
sino de la distancia entre lo que la red paga por recibir energía y lo que
cobra por entregarla. Es, por tanto, un margen de origen regulatorio y no
tecnológico, lo que condiciona su permanencia: una reforma tarifaria puede
estrecharlo sin que cambie nada en las instalaciones. Elaboración propia
desde las series de bolsa y de tarifa.}

La amplitud de esa banda es lo que hace viable el ejercicio, y su origen
es enteramente regulatorio: no procede de una diferencia de costos de
producción sino de la distancia entre lo que la red paga por recibir
energía y lo que cobra por entregarla. Entre la media de bolsa
(\uni{182.5}{COP/kWh}) y el costo unitario oficial, \uni{792.06}{}, hay un
factor superior a cuatro. Cualquier precio intermedio deja a las dos
partes mejor que la alternativa de operar contra la red.

\begin{trampabox}
A lo largo del documento aparecen varias cifras asociadas al «precio al
que la red compra» y conviene deshacer la confusión de una vez, porque
corresponden a papeles distintos de un mismo referente económico. El
\emph{piso de la dinámica} del juego es una constante de
\uni{280}{COP/kWh}, heredada de la calibración inicial, que acota
inferiormente la búsqueda del equilibrio. La \emph{serie horaria de
bolsa}, con media de \uni{182.5}{COP/kWh}, no interviene en esa dinámica:
es la que liquida los escenarios que la requieren y la exportación a red
del excedente no vendido. Y el \emph{valor nominal} de los barridos de
sensibilidad del Capítulo~\ref{sec:rob} es también \uni{280}{}, coherente
con el piso del juego. Piso de la dinámica, serie de liquidación y nominal
de barridos son tres papeles y no deben confundirse.
\end{trampabox}

\subsection{Lo que cuesta generar}
\label{sub:precios-lcoe}

Falta un tercer precio, que no se compra ni se vende: lo que le cuesta a
la propia institución el kilovatio-hora que produce. Para un sistema
fotovoltaico ya instalado, el costo marginal de producir una unidad
adicional es prácticamente nulo, de modo que el parámetro relevante es el
costo nivelado, es decir, la inversión amortizada sobre la energía que el
sistema entregará a lo largo de su vida.

El valor adoptado es de \uni{225}{COP/kWh} para las cuatro instituciones
con inversor de un fabricante y \uni{210}{} para Cesmag, que tiene otro.
Tras el ajuste por la irradiancia local de Pasto, los valores efectivos en
el simulador quedan en torno a \uni{241}{COP/kWh} para las primeras y
\uni{225}{} para Cesmag, lo que introduce una heterogeneidad del
\pct{6.67} entre ambas.

El rango se sustenta en tres referencias de distinto alcance: el costo
nivelado ponderado global de la energía solar fotovoltaica que reporta la
agencia internacional de energías renovables; los costos de referencia
para generación solar en Colombia del plan indicativo de expansión, entre
\num{300} y \uni{450}{COP/kWh} para escalas utilitaria y media; y el
análisis de incentivos económicos a la energía solar en el país, que sitúa
el costo nivelado efectivo de sistemas distribuidos en el rango de
\num{200} a \uni{250}{COP/kWh} bajo condiciones de financiación
favorables.

\begin{detallebox}
El valor adoptado es conservador dentro de ese rango, y esa elección tiene
una consecuencia contraintuitiva que el Capítulo~\ref{sec:rob} demuestra
de forma cuantitativa: el resultado del estudio \textbf{no depende de este
parámetro}. El análisis de sensibilidad global asigna al costo de generar
un índice de influencia sobre la ganancia del orden de $10^{-7}$, es decir,
indistinguible de cero. La razón es estructural: lo que el mercado transa
es excedente después de autoconsumo, y el volumen de ese excedente lo
determina el recurso disponible, no el costo contable de haberlo
producido. Se documenta aquí, no obstante, con el mismo detalle que el
resto, porque un parámetro que resulta irrelevante \emph{a posteriori} no
podía saberse irrelevante \emph{a priori}.
\end{detallebox}

\subsection{Una provisionalidad declarada}
\label{sub:precios-provisional}

El nivel de tensión asumido para las cinco instituciones, del cual
dependen todas las tarifas de este capítulo, es el nivel 2, y es
provisional hasta su verificación contra facturas reales. La publicación
tarifaria del comercializador respalda los valores de la categoría no
residencial en ese nivel y el factor de $1{,}20$ entre las dos categorías;
lo que queda abierto es únicamente la verificación del nivel de tensión
efectivo de cada institución, que la revisión regulatoria del proyecto
sitúa en el nivel 1.

Se declara aquí y no en un anexo porque afecta al valor absoluto de todas
las cifras económicas del documento. No afecta, en cambio, al
ordenamiento entre mecanismos, que es lo que el trabajo pretende
establecer: un cambio de nivel de tensión desplaza el techo de la banda
para las cinco instituciones por igual, y el Capítulo~\ref{sec:rob}
muestra que el ordenamiento sobrevive a variaciones de ese techo muy
superiores a la diferencia entre los dos niveles.

En síntesis, con este capítulo queda completo el insumo del modelo: a las
dos matrices de energía del Capítulo~\ref{sec:prep} se suman ahora la
serie horaria de bolsa, la matriz mensual de tarifa por institución y el
costo nivelado de generar, junto con la banda de \uni{182.5}{} a
\uni{792.06}{COP/kWh} dentro de la cual puede existir un intercambio que
convenga a las dos partes. El capítulo siguiente describe el mecanismo que
busca un precio dentro de esa banda.
